% Options for packages loaded elsewhere
\PassOptionsToPackage{unicode}{hyperref}
\PassOptionsToPackage{hyphens}{url}
\PassOptionsToPackage{dvipsnames,svgnames,x11names}{xcolor}
%
\documentclass[
  11pt,
]{article}

\usepackage{amsmath,amssymb}
\usepackage{iftex}
\ifPDFTeX
  \usepackage[T1]{fontenc}
  \usepackage[utf8]{inputenc}
  \usepackage{textcomp} % provide euro and other symbols
\else % if luatex or xetex
  \usepackage{unicode-math}
  \defaultfontfeatures{Scale=MatchLowercase}
  \defaultfontfeatures[\rmfamily]{Ligatures=TeX,Scale=1}
\fi
\usepackage{lmodern}
\ifPDFTeX\else  
    % xetex/luatex font selection
\fi
% Use upquote if available, for straight quotes in verbatim environments
\IfFileExists{upquote.sty}{\usepackage{upquote}}{}
\IfFileExists{microtype.sty}{% use microtype if available
  \usepackage[]{microtype}
  \UseMicrotypeSet[protrusion]{basicmath} % disable protrusion for tt fonts
}{}
\makeatletter
\@ifundefined{KOMAClassName}{% if non-KOMA class
  \IfFileExists{parskip.sty}{%
    \usepackage{parskip}
  }{% else
    \setlength{\parindent}{0pt}
    \setlength{\parskip}{6pt plus 2pt minus 1pt}}
}{% if KOMA class
  \KOMAoptions{parskip=half}}
\makeatother
\usepackage{xcolor}
\usepackage[margin=1in]{geometry}
\setlength{\emergencystretch}{3em} % prevent overfull lines
\setcounter{secnumdepth}{-\maxdimen} % remove section numbering
% Make \paragraph and \subparagraph free-standing
\ifx\paragraph\undefined\else
  \let\oldparagraph\paragraph
  \renewcommand{\paragraph}[1]{\oldparagraph{#1}\mbox{}}
\fi
\ifx\subparagraph\undefined\else
  \let\oldsubparagraph\subparagraph
  \renewcommand{\subparagraph}[1]{\oldsubparagraph{#1}\mbox{}}
\fi


\providecommand{\tightlist}{%
  \setlength{\itemsep}{0pt}\setlength{\parskip}{0pt}}\usepackage{longtable,booktabs,array}
\usepackage{calc} % for calculating minipage widths
% Correct order of tables after \paragraph or \subparagraph
\usepackage{etoolbox}
\makeatletter
\patchcmd\longtable{\par}{\if@noskipsec\mbox{}\fi\par}{}{}
\makeatother
% Allow footnotes in longtable head/foot
\IfFileExists{footnotehyper.sty}{\usepackage{footnotehyper}}{\usepackage{footnote}}
\makesavenoteenv{longtable}
\usepackage{graphicx}
\makeatletter
\def\maxwidth{\ifdim\Gin@nat@width>\linewidth\linewidth\else\Gin@nat@width\fi}
\def\maxheight{\ifdim\Gin@nat@height>\textheight\textheight\else\Gin@nat@height\fi}
\makeatother
% Scale images if necessary, so that they will not overflow the page
% margins by default, and it is still possible to overwrite the defaults
% using explicit options in \includegraphics[width, height, ...]{}
\setkeys{Gin}{width=\maxwidth,height=\maxheight,keepaspectratio}
% Set default figure placement to htbp
\makeatletter
\def\fps@figure{htbp}
\makeatother

\usepackage{array}
\usepackage{booktabs}
\usepackage{graphicx}
\usepackage{longtable}
\usepackage{needspace}
\usepackage{caption}
\usepackage{float}
\usepackage{fancyhdr}
\pagestyle{fancy}
\fancyhf{}
\fancyhead[L]{Video Quality Analysis Report}
\fancyhead[R]{Fifth Avenue Studios}
\fancyfoot[L]{\today}
\fancyfoot[C]{\thepage}
\fancyfoot[R]{Gil Raitses}
\setlength{\extrarowheight}{1pt}
\renewcommand{\arraystretch}{1.0}
\makeatletter
\makeatother
\makeatletter
\makeatother
\makeatletter
\@ifpackageloaded{caption}{}{\usepackage{caption}}
\AtBeginDocument{%
\ifdefined\contentsname
  \renewcommand*\contentsname{Table of contents}
\else
  \newcommand\contentsname{Table of contents}
\fi
\ifdefined\listfigurename
  \renewcommand*\listfigurename{List of Figures}
\else
  \newcommand\listfigurename{List of Figures}
\fi
\ifdefined\listtablename
  \renewcommand*\listtablename{List of Tables}
\else
  \newcommand\listtablename{List of Tables}
\fi
\ifdefined\figurename
  \renewcommand*\figurename{Figure}
\else
  \newcommand\figurename{Figure}
\fi
\ifdefined\tablename
  \renewcommand*\tablename{Table}
\else
  \newcommand\tablename{Table}
\fi
}
\@ifpackageloaded{float}{}{\usepackage{float}}
\floatstyle{ruled}
\@ifundefined{c@chapter}{\newfloat{codelisting}{h}{lop}}{\newfloat{codelisting}{h}{lop}[chapter]}
\floatname{codelisting}{Listing}
\newcommand*\listoflistings{\listof{codelisting}{List of Listings}}
\makeatother
\makeatletter
\@ifpackageloaded{caption}{}{\usepackage{caption}}
\@ifpackageloaded{subcaption}{}{\usepackage{subcaption}}
\makeatother
\makeatletter
\@ifpackageloaded{tcolorbox}{}{\usepackage[skins,breakable]{tcolorbox}}
\makeatother
\makeatletter
\@ifundefined{shadecolor}{\definecolor{shadecolor}{rgb}{.97, .97, .97}}
\makeatother
\makeatletter
\makeatother
\makeatletter
\makeatother
\ifLuaTeX
  \usepackage{selnolig}  % disable illegal ligatures
\fi
\IfFileExists{bookmark.sty}{\usepackage{bookmark}}{\usepackage{hyperref}}
\IfFileExists{xurl.sty}{\usepackage{xurl}}{} % add URL line breaks if available
\urlstyle{same} % disable monospaced font for URLs
\hypersetup{
  pdftitle={Video Quality Analysis Report},
  colorlinks=true,
  linkcolor={blue},
  filecolor={Maroon},
  citecolor={Blue},
  urlcolor={Blue},
  pdfcreator={LaTeX via pandoc}}

\title{Video Quality Analysis Report}
\usepackage{etoolbox}
\makeatletter
\providecommand{\subtitle}[1]{% add subtitle to \maketitle
  \apptocmd{\@title}{\par {\large #1 \par}}{}{}
}
\makeatother
\subtitle{Flaco Owl Documentary Footage - Fifth Avenue Studios}
\author{}
\date{2025-11-13}

\begin{document}
\maketitle
\begin{abstract}
This report analyzes 24.70 seconds of video footage documenting Flaco
the owl's first landing on Fifth Avenue on February 4, 2024, the night
he was freed from the Central Park Zoo. This clip distinguishes itself
as the longest of the three clips (24.70s) and features the best
sharpness score (66.03) among all clips, making it the sharpest footage
available. Technical analysis reveals HD resolution (1280 × 720 pixels,
30 fps) with overall quality score of 49.4/100, requiring
post-processing for brightness adjustment and stability. Unlike the
other clips, this footage contains one clear/focused segment
(13.5-16.0s, 2.5 seconds) with a quality score of 65.0, providing a
ready-to-use segment without post-processing. The footage achieves
98.0\% usability (24.21 seconds usable) with 16.0\% meeting high-quality
standards (3.95 seconds), and uniquely offers two separate usable
segments (0-19s and 20-24.7s) providing editing flexibility. Content
analysis shows 1 frame with dog detection, the only clip with any animal
detection. Valuation analysis indicates a rate of \$105.47 per usable
second, resulting in an estimated value range of \$2,553 - \$2,761. This
clip offers the longest duration, best sharpness, and includes a
production-ready clear segment.
\end{abstract}
\ifdefined\Shaded\renewenvironment{Shaded}{\begin{tcolorbox}[borderline west={3pt}{0pt}{shadecolor}, frame hidden, interior hidden, sharp corners, breakable, boxrule=0pt, enhanced]}{\end{tcolorbox}}\fi

\needspace{8cm}
\section{Executive Summary}

This report provides a technical and narrative analysis of video footage
documenting Flaco the owl's first-night landing on Fifth Avenue after
being freed from the Central Park Zoo. The footage has been evaluated
for documentary production suitability and content significance.

\begin{table}[h]
\centering
\small
\caption{Executive Summary Findings}
\label{tab-exec-summary}
\begin{tabular}{p{3cm}p{8cm}}
\toprule
\textbf{Finding} & \textbf{Value} \\
\midrule
Total Duration & 24.70 seconds \\
Usable Footage & 24.21 seconds (98.0\%) \\
High Quality Segments & 3.95 seconds (16.0\%) \\
Narrative Significance & 90/100 (high importance) \\
Estimated Value Range & \$2,553 - \$2,761 \\
\bottomrule
\end{tabular}
\end{table}

\begin{figure}[h]
\centering
\includegraphics[width=0.85\textwidth]{figures_3789/value_range.pdf}
\caption{Longest clip duration (24.70s) provides the most footage despite lower per-second rate, resulting in competitive total value range (\$2,553 - \$2,761). Bar chart comparing Conservative (\$2,553), Recommended and Optimistic (\$2,761) valuation estimates.}
\label{fig-value-range}
\end{figure}

\needspace{8cm}
\section{Technical Specifications}

\begin{table}[h]
\centering
\small
\caption{Technical Specifications}
\label{tab-tech-specs}
\begin{tabular}{p{3cm}p{7.5cm}}
\toprule
\textbf{Property} & \textbf{Value} \\
\midrule
Resolution & 1280 × 720 pixels (HD 720p, Landscape) \\
Frame Rate & 30.00 fps \\
Duration & 24.70 seconds \\
Total Frames & 741 frames \\
File Size & 7.25 MB \\
File Format & QuickTime (.mov) \\
Video Codec & H.264 (AVC) \\
Color Depth & 8-bit \\
Pixel Format & YUV 4:2:0 \\
Aspect Ratio & 16:9 (Landscape) \\
\bottomrule
\end{tabular}
\end{table}

\needspace{10cm}
\section{Visual Quality Metrics}

\subsection{Core Quality Metrics}

\begin{table}[h]
\centering
\small
\caption{Core Quality Metrics}
\label{tab-quality-metrics}
\begin{tabular}{p{3.5cm}p{2cm}p{5.5cm}}
\toprule
\textbf{Metric} & \textbf{Score} & \textbf{Assessment} \\
\midrule
Sharpness & 66.03 & Moderate Sharpness \\
Brightness & 59.7/255 & Underexposed \\
Contrast & 43.43 & High Contrast \\
Noise Level & 7.92 & Clean \\
Exposure & 92.6/100 & Good Exposure \\
Stability & 6.61 & Unstable/Shaky \\
Smoothness & 0.934 & Very Smooth \\
Color Compression & 0.979 & Minimal Compression \\
\bottomrule
\end{tabular}
\end{table}

\needspace{12cm}
\subsection{Quality Assessment}

\textbf{Strengths:}

\begin{table}[h]
\centering
\small
\caption{Quality Strengths}
\label{tab-strengths}
\begin{tabular}{p{2cm}p{7cm}}
\toprule
\textbf{Strength} & \textbf{Description} \\
\midrule
Sharpness & Moderate sharpness (66.03) better than other clips \\
Exposure & Well-exposed footage with minimal over/under exposure \\
Noise Level & Clean footage with minimal noise (7.92) \\
Smoothness & High temporal smoothness (0.934) \\
Compression & Minimal color compression artifacts (0.979) \\
\bottomrule
\end{tabular}
\end{table}

\textbf{Post-Processing Considerations:}

\begin{table}[h]
\centering
\small
\caption{Post-Processing Requirements}
\label{tab-post-processing}
\begin{tabular}{p{2cm}p{4cm}p{3cm}}
\toprule
\textbf{Issue} & \textbf{Description} & \textbf{Required Action} \\
\midrule
Brightness & Footage is underexposed (brightness: 59.7) & Requires brightness adjustment \\
Stability & Camera shake present (stability: 6.61) & Requires stabilization \\
\bottomrule
\end{tabular}
\end{table}

\begin{figure}[h]
\centering
\includegraphics[width=0.7\textwidth]{figures_3789/quality_metrics_radar.pdf}
\caption{This clip features the best sharpness score (66.03) among all three clips, making it the sharpest footage available. Lower scores in brightness and stability indicate areas requiring post-processing enhancement. Radar chart showing normalized quality scores across six dimensions.}
\label{fig-quality-radar}
\end{figure}

\needspace{22cm}
\section{Content Significance Analysis}

\subsection{Subject Detection}

\begin{table}[h]
\centering
\small
\caption{Subject Detection Metrics}
\label{tab-subject-detection}
\begin{tabular}{p{4.5cm}p{5.5cm}}
\toprule
\textbf{Metric} & \textbf{Value} \\
\midrule
Frames Analyzed & 50 frames \\
Frames with Subject (Bird/Owl) & 0 frames (0.0\%) \\
Frames with Dog & 1 frame \\
Frames with Interaction & 0 frames \\
Subject Visible Throughout & No \\
Interaction Moments Detected & No \\
Rare/Unique Event & Yes \\
\bottomrule
\end{tabular}
\end{table}

\subsection{Event Significance}

This footage documents a rare and historic event: Flaco the owl's first
landing on Fifth Avenue on the night he was freed from the Central Park
Zoo. This represents an important moment in the narrative arc of the
documentary, capturing the first public appearance of Flaco after his
escape. This clip is unique among the three as the only one with any
animal detection (1 frame with dog), and provides the longest continuous
footage (24.70 seconds) for extended narrative sequences.

\needspace{20cm}
\section{Narrative Significance (Story Value)}

\needspace{12cm}
\subsection{Narrative Score: 90/100}

\begin{table}[h]
\centering
\small
\caption{Narrative Significance Factors}
\label{tab-narrative-factors}
\begin{tabular}{p{3cm}p{6cm}}
\toprule
\textbf{Factor} & \textbf{Contribution} \\
\midrule
Rare/Unique Event & First-night landing on Fifth Ave \\
Historic Moment & First appearance after zoo escape \\
Subject Visibility & Subject not detected in analyzed frames \\
Multiple-Use Potential & Can be used multiple times throughout documentary \\
Story Importance & High importance \\
\bottomrule
\end{tabular}
\end{table}

\textbf{Multiple-Use Potential:} This footage can be strategically
reused throughout the documentary in:

\begin{table}[h]
\centering
\small
\caption{Multiple-Use Applications}
\label{tab-multiple-use}
\begin{tabular}{p{3.5cm}p{7.5cm}}
\toprule
\textbf{Use Case} & \textbf{Description} \\
\midrule
Opening Sequence & Can be used in documentary opening \\
Transitional Moments & Suitable for scene transitions \\
Key Narrative Scenes & Can feature in important story moments \\
Closing Sequences & Appropriate for documentary conclusion \\
\bottomrule
\end{tabular}
\end{table}

\begin{figure}[h]
\centering
\includegraphics[width=0.7\textwidth]{figures_3789/narrative_significance.pdf}
\caption{Narrative score of 90/100 reflects historic moment value enhanced by longest duration (24.70s) providing extended narrative coverage and unique animal detection (dog). Pie chart showing narrative significance score of 90/100.}
\label{fig-narrative}
\end{figure}

\needspace{25cm}
\section{Usability Analysis}

\subsection{Usability Breakdown}

\begin{figure}[h]
\centering
\includegraphics[width=0.7\textwidth]{figures_3789/usability_breakdown.pdf}
\caption{Longest clip (24.70s) with 98.0\% usability, uniquely offering two separate usable segments providing editing flexibility. Bar chart comparing total duration (24.70s, 100\%), usable seconds (24.21s, 98.0\%) and high-quality seconds (3.95s, 16.0\%).}
\label{fig-usability-breakdown}
\end{figure}

\subsection{Per-Second Breakdown}

\begin{table}[h]
\centering
\small
\caption{Per-Second Breakdown}
\label{tab-per-second}
\begin{tabular}{p{2.5cm}cc}
\toprule
\textbf{Category} & \textbf{Seconds} & \textbf{Percentage} \\
\midrule
Total Duration & 24.70 & 100.0\% \\
Usable Seconds & 24.21 & 98.0\% \\
High Quality Seconds & 3.95 & 16.0\% \\
\bottomrule
\end{tabular}
\end{table}

\textbf{Usable Segments:} 2 segments (0.00s - 19.00s, 20.00s - 24.70s)

\subsection{Clear, Focused Segments}

One clear/focused segment identified: 13.50s - 16.00s (2.50 seconds)
with quality score of 65.0. The footage requires post-processing
(brightness adjustment and stabilization) to achieve optimal clarity for
broadcast use.

\needspace{25cm}
\section{Valuation Analysis}

\subsection{Valuation Breakdown}

\begin{figure}[h]
\centering
\includegraphics[width=0.75\textwidth]{figures_3789/valuation_breakdown.pdf}
\caption{Lower per-second rate (\$105.47) reflects HD 720p resolution multiplier (0.75×) and quality multiplier (0.50×), but longest duration compensates for lower rate. Waterfall chart showing how base rate (\$50.00) is modified by sequential multipliers: Quality (0.50×), Resolution (0.75×), Content Significance (1.50×), Narrative Significance (2.50×) and Multiple-Use (1.50×).}
\label{fig-valuation-breakdown}
\end{figure}

\needspace{8cm}
\subsection{Rate Calculation}

\enlargethispage{2cm}
\vspace{-0.5cm}
\begin{table}[H]
\centering
\footnotesize
\caption{Rate Calculation}
\label{tab-rate-calculation}
\begin{tabular}{p{3.5cm}p{1.8cm}p{4.2cm}}
\toprule
\textbf{Factor} & \textbf{Multiplier} & \textbf{Value} \\
\midrule
Base Rate (per usable second) & 1.00× & \$50.00 \\
Quality Multiplier & 0.50× & \\
Resolution Multiplier & 0.75× & \\
Content Significance Multiplier & 1.50× & \\
Narrative Significance Multiplier & 2.50× & \\
Multiple-Use Multiplier & 1.50× & \\
\textbf{Final Rate per Second} & & \textbf{\$105.47} \\
\bottomrule
\end{tabular}
\end{table}

\needspace{8cm}
\subsection{Estimated Value Range}

\begin{table}[h]
\centering
\small
\caption{Estimated Value Range}
\label{tab-value-range}
\begin{tabular}{p{3cm}p{5cm}p{3cm}}
\toprule
\textbf{Estimate Type} & \textbf{Calculation} & \textbf{Value} \\
\midrule
Conservative & All usable seconds × \$105.47 & \$2,553 \\
Optimistic & Premium segments at 1.5× rate & \$2,761 \\
\textbf{Recommended Range} & & \textbf{\$2,553 - \$2,761} \\
\bottomrule
\end{tabular}
\end{table}

\needspace{12cm}
\subsection{Value Factors}

\begin{table}[h]
\centering
\small
\caption{Value Factors}
\label{tab-value-factors}
\begin{tabular}{p{3cm}p{7cm}}
\toprule
\textbf{Factor} & \textbf{Description} \\
\midrule
Usability & High usability (98.0\% usable) - most of the clip is production-ready \\
Archival Value & Rare/unique event adds archival and narrative value \\
Story Value & Story moment with reuse potential increases value \\
Multiple-Use & Footage can be used multiple times throughout documentary structure \\
Editing Flexibility & Multiple usable segments (2) provide editing flexibility \\
\bottomrule
\end{tabular}
\end{table}

\section{Conclusion}

This footage represents an important moment in Flaco the owl's story,
documenting his first landing on Fifth Avenue after being freed from the
Central Park Zoo. This clip distinguishes itself as the longest of the
three (24.70 seconds) and features the best sharpness score (66.03)
among all clips, making it the sharpest footage available. Unlike the
other clips, this footage contains one clear/focused segment
(13.5-16.0s, 2.5 seconds) with a quality score of 65.0, providing a
production-ready segment that requires no post-processing. The footage
uniquely offers two separate usable segments (0-19s and 20-24.7s)
providing editing flexibility, and is the only clip with any animal
detection (1 frame with dog). While brightness and stability require
post-processing for the full clip, the combination of longest duration,
best sharpness, and a ready-to-use clear segment makes this footage
valuable for extended narrative sequences and editing flexibility. The
estimated licensing value range of \$2,553 - \$2,761 reflects the longer
duration compensating for the lower per-second rate due to HD 720p
resolution.



\end{document}
